\documentclass[11pt]{article}
\usepackage{amsmath,amssymb,amsthm,mathrsfs}
\usepackage[margin=1in]{geometry}
\usepackage{hyperref}

\newtheorem{theorem}{Theorem}
\newtheorem{proposition}[theorem]{Proposition}
\newtheorem{lemma}[theorem]{Lemma}
\newtheorem{corollary}[theorem]{Corollary}
\theoremstyle{definition}
\newtheorem{remark}[theorem]{Remark}
\newtheorem{convention}[theorem]{Convention}
\newtheorem{question}[theorem]{Question}

\DeclareMathOperator{\ordop}{ord}
\newcommand{\hgt}{\operatorname{ht}}   % NOT \ht -- collides with a TeX primitive
\newcommand{\Z}{\mathbb{Z}}
\newcommand{\Q}{\mathbb{Q}}
\newcommand{\Lam}{\Lambda}
\newcommand{\F}{\mathcal{F}}
\newcommand{\Mult}[1]{M_{#1}}

\title{The cross-rank obstruction for ribbon operators does not need\\
       a single parameter: free height weights force Murnaghan--Nakayama}
\author{Clio}
\date{11 September 2026 (cycle 2)}

\begin{document}
\maketitle

\begin{abstract}
For $e\ge1$ let $R_e(t)$ be the operator on the ring of symmetric functions that adds a
connected $e$-ribbon with weight $t^{\hgt}$.  I proved previously that for $e\ne f$ the
operators $R_e(t)$ and $R_f(t)$ commute only at $t=-1$.  That theorem was proved for a
\emph{single} parameter $t$, and the height weight $t^{N}$ was assumed multiplicative in
$N$.  Following LLT's 1995 move --- replace a parameter raised to a statistic by a
parameter indexed by that statistic --- I remove the assumption entirely: give each
height value $N$ its own weight $w_N$, with no relation between them, and allow the
weight functions of $R_e$ and $R_f$ to be \emph{unrelated to each other} as well.

\textbf{Result.}  The obstruction survives, and it survives in the strongest form the
question admits.  For $1\le e<f$ and arbitrary weight functions
$w:\{0,\dots,e-1\}\to K$, $\bar w:\{0,\dots,f-1\}\to K$ over an integral domain,
$[R_e^{w},R_f^{\bar w}]=0$ if and only if one of the operators is zero or
\[
   w_N=(-1)^N w_0\ \ (N\le e-1),\qquad \bar w_N=(-1)^N\bar w_0\ \ (N\le f-1),
\]
i.e.\ both are scalar multiples of multiplication by a power sum.  \emph{Every parameter
the pair can see is pinned by the pair.}  In the coordinates of the brief
($\tau_N=(-1)^Nw_N$, so that single-$t$ is $\tau_N=(-t)^N$) the commuting locus is the
single point $\tau\equiv1$.  This is outcome~1: the single-$t$ hypothesis of the forcing
theorem was never load-bearing, and the theorem strengthens to arbitrary height weights.

The proof isolates \emph{why}.  The commutator splits into a one-bead and a two-bead
sector, and the two sectors cut out different loci: the two-bead sector alone forces
both weights to be geometric with reciprocal ratios ($w_N=s^{N}$, $\bar w_N=s^{-N}$),
which is exactly the restoration of a single parameter; the one-bead sector alone forces
the anchor outright.  The mechanism in the one-bead sector is a dichotomy internal to the
height statistic: an interior site of a long ribbon is counted by the height on the route
that steps over it and not counted on the route that factors the ribbon there.  The two
readings differ by exactly one, so closure forces $w_{N+1}=-w_{N}$, and the familiar
scalar $1+t$ is only the instance $N=0$.
\end{abstract}

\section{Conventions, and the two things that had to be pinned first}

\begin{convention}\label{conv:main}
$K$ is an integral domain; $\Lam_K$ is the ring of symmetric functions over $K$ with
Schur basis $\{s_\lambda\}$ and Hall inner product $\langle s_\lambda,s_\mu\rangle=
\delta_{\lambda\mu}$.  For a partition $\lambda$, $M(\lambda)=\{\lambda_j-j:j\ge1\}$ is
its Maya set (beads); $m(x)=[x\in M]$.  For $x<y$ we write
$\#(M\cap(x,y))$ for the number of beads strictly between $x$ and $y$.
For a connected border strip (ribbon) $\mu/\lambda$, $\hgt(\mu/\lambda)$ is its number of
rows minus one.
\end{convention}

\begin{convention}[the deformed operator]\label{conv:def}
Fix $e\ge1$ and a \emph{weight function} $w:\{0,1,\dots,e-1\}\to K$, written $w_N$.
Define $R_e^{w}:\Lam_K\to\Lam_K$ by
\[
   R_e^{w}\,s_\lambda\;=\;\sum_{\mu}\;w_{\hgt(\mu/\lambda)}\;s_\mu ,
\]
the sum over $\mu\supset\lambda$ with $\mu/\lambda$ a connected $e$-ribbon.  (The height
of an $e$-ribbon lies in $\{0,\dots,e-1\}$, so $R_e^w$ sees exactly these $e$ values and
no more.)  The classical operator is $w_N=t^{N}$:  $R_e^{(t^\bullet)}=R_e(t)$.
\end{convention}

Two conventions had to be settled before any deformation, because getting either wrong
silently answers a different question.

\subsection*{First: the sign is not a second parameter}

The fermionic normal form of \cite{Clio0906} reads
$R_e(t)=\sum_{b}\psi_{b+e}\psi_b^{*}(-t)^{N_{(b,b+e)}}$ with
$N_{(b,b+e)}=\sum_{b<j<b+e}n_j$, and the deformation proposed in the brief is
$R_e^{\tau}=\sum_b\psi_{b+e}\psi_b^{*}\tau_{N_{(b,b+e)}}$.  The $(-1)^N$ inside $(-t)^N$
is \emph{not} part of the weight: it cancels the Jordan--Wigner sign.  Explicitly
\cite[Lem.~5]{Clio0906},
$\psi_{b+e}\psi_b^{*}|M\rangle=(-1)^{\#(M\cap(b,b+e))}|M'\rangle$, so
\[
   R_e^{\tau}\,|\lambda\rangle
   \;=\;\sum_{\mu}(-1)^{\hgt(\mu/\lambda)}\tau_{\hgt(\mu/\lambda)}\,|\mu\rangle ,
   \qquad\text{that is}\qquad
   \boxed{\,w_N=(-1)^N\tau_N\,}.
\]
The map $\tau\mapsto w$ is a diagonal linear isomorphism of the parameter space, so a
vanishing locus computed in one coordinate transports verbatim to the other; there is no
double counting and nothing extra to deform.  Under it:
\[
 \text{single-}t:\ \tau_N=(-t)^N\ \Longleftrightarrow\ w_N=t^N;\qquad
 \text{anchor:}\ \tau\equiv1\ \Longleftrightarrow\ w_N=(-1)^N\ \Longleftrightarrow\ t=-1 .
\]
\textbf{I work throughout on the partition side, in the coordinates $w$, where no fermion
sign occurs at all}, and translate to $\tau$ only in Corollary~\ref{cor:tau}.  This makes
the hazard structurally impossible rather than merely avoided.

\subsection*{Second: shared versus independent weights}

LLT's move gives one parameter per height value, \emph{shared} across all ribbon sizes:
$w^{(e)}=w^{(f)}$.  The outer bound gives each $R_e$ its own weight function.  I prove
the theorem for \textbf{independent} weight functions; the shared statement is the
special case.  Nothing below assumes any relation between $w$ and $\bar w$.

\subsection*{What is being tested}

\begin{question}\label{q:main}
For which pairs $(w,\bar w)$ does $[R_e^{w},R_f^{\bar w}]=0$, $e\ne f$?
\end{question}

The point of the question is that $R_e^w$ has $e-1$ free parameters after normalisation,
so $e\le2$ cannot distinguish a free weight from a single $t$: $R_1^w$ has none and
$R_2^w$ has one, exactly as many as $t$.  The first genuinely new freedom is
$w_2$ versus $w_1^2$ at $e=3$; accordingly $[R_1,R_2]$ is a degenerate test and is not
used anywhere below.

\section{Realisability: every local bead pattern occurs}

Everything is decided inside one window of length $e+f$.  The following says the window
may be prescribed arbitrarily.

\begin{lemma}\label{lem:realise}
Let $g\ge2$ and let $S\subseteq\{1,\dots,g-1\}$.  Then there is a partition $\lambda$
with $0\in M(\lambda)$, $g\notin M(\lambda)$ and $M(\lambda)\cap(0,g)=S$.
\end{lemma}

\begin{proof}
Put $r=1+|S|$ and $M=\Z_{\le -r-1}\cup\{0\}\cup S$.  Then $M$ contains all sufficiently
negative integers and is bounded above, and its charge is
$\#(M\cap\Z_{\ge0})-\#(\Z_{<0}\setminus M)=(1+|S|)-r=0$; hence $M=M(\lambda)$ for a
unique partition $\lambda$.  By construction $0\in M$, $g\notin M$ (as
$S\subseteq\{1,\dots,g-1\}$) and $M\cap(0,g)=S$.  The $r$ deleted sites lie below $0$ and
so do not meet the window.
\end{proof}

\section{The one-bead matrix element with free weights}

Fix $1\le e<f$ and put $g=e+f$.  Throughout, $M=M(\lambda)$ is a Maya set with $b\in M$
and $b+g\notin M$, and $\mu$ is the partition with
$M(\mu)=(M\setminus\{b\})\cup\{b+g\}$.  Write
\[
  m_e=m(b+e),\quad m_f=m(b+f),
\]
\[
  P_e=\#(M\cap(b,b+e)),\quad S_e=\#(M\cap(b+e,b+g)),\quad
  P_f=\#(M\cap(b,b+f)),\quad S_f=\#(M\cap(b+f,b+g)).
\]

\begin{lemma}[one-bead element]\label{lem:onebead}
With the notation above,
\[
 \boxed{\ \bigl\langle s_\mu,\ [R_e^{w},R_f^{\bar w}]\,s_\lambda\bigr\rangle
   \;=\;(1-2m_f)\,w_{S_f}\,\bar w_{P_f}\;-\;(1-2m_e)\,w_{P_e}\,\bar w_{S_e}. \ }
\]
All four subscripts are legal: $P_e,S_f\le e-1$ and $P_f,S_e\le f-1$.
\end{lemma}

\begin{proof}
By Convention~\ref{conv:def} and \cite[Lem.~4]{Clio0906}, $R_g^{\bullet}$ acts on Maya
sets by legal bead moves ($b'\in M$, $b'+g\notin M$), with the height of the resulting
ribbon equal to the number of beads strictly jumped.  Hence
\[
   \bigl\langle s_\mu,\ R_f^{\bar w}R_e^{w}s_\lambda\bigr\rangle
   \;=\;\sum_{\text{routes}} w_{h_1}\bar w_{h_2},
\]
over pairs of legal moves $M\xrightarrow{(x,e)}M_1\xrightarrow{(y,f)}M(\mu)$, with
$h_1=\#(M\cap(x,x+e))$ and $h_2=\#(M_1\cap(y,y+f))$.

\emph{Which routes exist.}  We have
$M(\mu)=\bigl((M\setminus\{x\})\cup\{x+e\}\bigr)\setminus\{y\}\cup\{y+f\}$.  If no
removed site coincides with an added one then $|M\,\triangle\,M(\mu)|=4$; but
$|M\,\triangle\,M(\mu)|=2$ here.  Since $x\ne x+e$ and $y\ne y+f$, the only
possibilities are $y=x+e$ or $x=y+f$.

\emph{Route B ($y=x+e$).}  Then $M(\mu)=(M\setminus\{x\})\cup\{x+g\}$, so $x=b$,
$y=b+e$.  Legality of the first move needs $b\in M$ (given) and $b+e\notin M$, i.e.\
$m_e=0$; the second move is then automatically legal ($b+e\in M_1$, and
$b+g\notin M_1$ because $b+g\notin M$ and $M_1$ added only $b+e$).  Its heights:
$h_1=\#(M\cap(b,b+e))=P_e$, and since $M_1$ differs from $M$ only at $b$ and $b+e$, both
outside the \emph{open} interval $(b+e,b+g)$, $h_2=\#(M\cap(b+e,b+g))=S_e$.  Contribution
$w_{P_e}\bar w_{S_e}$, present iff $m_e=0$.

\emph{Route C ($x=y+f$).}  Then $M(\mu)=(M\setminus\{y\})\cup\{y+g\}$, so $y=b$,
$x=b+f$.  Legality of the first move needs $b+f\in M$, i.e.\ $m_f=1$, and
$b+g\notin M$ (given); the second is then legal.  Heights:
$h_1=\#(M\cap(b+f,b+g))=S_f$, and $M_1$ differs from $M$ only at $b+f$ and $b+g$, both
outside $(b,b+f)$, so $h_2=\#(M\cap(b,b+f))=P_f$.  Contribution
$w_{S_f}\bar w_{P_f}$, present iff $m_f=1$.

Therefore
$\langle s_\mu,R_f^{\bar w}R_e^{w}s_\lambda\rangle
 =(1-m_e)\,w_{P_e}\bar w_{S_e}+m_f\,w_{S_f}\bar w_{P_f}$.
Running the same enumeration for the other order --- now the $f$-move is performed first
--- Route~B$'$ ($f$-move at $b$, then $e$-move at $b+f$) exists iff $m_f=0$ and
contributes $\bar w_{P_f}w_{S_f}$; Route~C$'$ ($f$-move at $b+e$, then $e$-move at $b$)
exists iff $m_e=1$ and contributes $\bar w_{S_e}w_{P_e}$.  Hence
$\langle s_\mu,R_e^{w}R_f^{\bar w}s_\lambda\rangle
 =(1-m_f)\,w_{S_f}\bar w_{P_f}+m_e\,w_{P_e}\bar w_{S_e}$, and subtracting gives the
boxed formula.  The subscript bounds hold because $P_e,S_f$ count beads in open intervals
of length $e$ and $P_f,S_e$ in open intervals of length $f$.
\end{proof}

\begin{remark}[what the formula says, in words]\label{rem:hinge}
The two monomials are the two ways of \emph{factoring} the long ribbon $b\to b+g$ at an
interior site: through the hinge $b+e$, or through the hinge $b+f$.  Remarkably, the pair
of heights $(P_e,S_e)$ of the factorisation at $b+e$ is the same whether that site is
occupied or empty --- what the occupancy changes is \emph{which order realises it}: if
$b+e$ is empty the $e$-move goes first, and if $b+e$ is occupied its resident bead must
vacate first, so the $f$-move goes first.  The monomial therefore appears on one side of
the commutator or the other according to $m_e$, which is the factor $(1-2m_e)$.  This is
the whole mechanism, and the next section is bookkeeping.
\end{remark}

Specialising $w=\bar w=(t^\bullet)$ recovers
$\bigl(m_e-m_f\bigr)(1+t)t^{N-1}$ with $N=\#(M\cap(b,b+g))$, since
$P_e+S_e=N-m_e$ and $P_f+S_f=N-m_f$: this is \cite[Thm.~5(1)]{Clio0907}.

\section{The forcing theorem}

\begin{theorem}[free height weights force Murnaghan--Nakayama]\label{thm:main}
Let $1\le e<f$, let $K$ be an integral domain, and let $w:\{0,\dots,e-1\}\to K$ and
$\bar w:\{0,\dots,f-1\}\to K$ be arbitrary.  Then
$\bigl[R_e^{w},R_f^{\bar w}\bigr]=0$ if and only if one of the following holds:
\begin{enumerate}
\item[(a)] $w\equiv0$ (equivalently $R_e^{w}=0$);
\item[(b)] $\bar w\equiv0$ (equivalently $R_f^{\bar w}=0$);
\item[(c)] $w_N=(-1)^{N}w_0$ for $0\le N\le e-1$ \emph{and}
           $\bar w_N=(-1)^{N}\bar w_0$ for $0\le N\le f-1$.
\end{enumerate}
In case (c), $R_e^{w}=w_0\,\Mult{p_e}$ and $R_f^{\bar w}=\bar w_0\,\Mult{p_f}$,
multiplication by power sums.
\end{theorem}

Only one of the four occupancy patterns of Lemma~\ref{lem:onebead} is used, namely
$m_e=1$, $m_f=0$; and only the one-bead sector is used.  Write it out.  By
Lemma~\ref{lem:realise} we may choose the window freely, so let
\[
  u=\#\bigl(M\cap(b,b+e)\bigr),\quad v=\#\bigl(M\cap(b+e,b+f)\bigr),\quad
  z=\#\bigl(M\cap(b+f,b+g)\bigr),
\]
which range independently over $0\le u\le e-1$, $0\le v\le f-e-1$, $0\le z\le e-1$
(the three open intervals have $e-1$, $f-e-1$ and $g-f-1=e-1$ sites).  With $m_e=1$,
$m_f=0$ we get $P_e=u$, $S_e=v+z$, $P_f=u+1+v$, $S_f=z$, and Lemma~\ref{lem:onebead}
gives the element $w_z\bar w_{u+v+1}+w_u\bar w_{v+z}$.  So:

\begin{lemma}[the forcing relation]\label{lem:B}
If $[R_e^{w},R_f^{\bar w}]=0$ then
\begin{equation}\label{eq:B}
   w_z\,\bar w_{u+v+1}\;+\;w_u\,\bar w_{v+z}\;=\;0
   \qquad\text{for all } 0\le u,z\le e-1,\ \ 0\le v\le f-e-1 .
\end{equation}
\end{lemma}

\begin{proof}
Immediate from Lemmas~\ref{lem:realise} and~\ref{lem:onebead}: given $(u,v,z)$, choose
$S\subseteq\{1,\dots,g-1\}$ with $e\in S$, $f\notin S$, $|S\cap(0,e)|=u$,
$|S\cap(e,f)|=v$, $|S\cap(f,g)|=z$, and take $\lambda$ from Lemma~\ref{lem:realise} and
$b=0$.
\end{proof}

\begin{proof}[Proof of Theorem~\ref{thm:main}]
\emph{Sufficiency.}  If $w\equiv0$ or $\bar w\equiv0$ one of the operators is zero and
the commutator vanishes.  Suppose (c).  Both operators are $K$-linear in their weights,
so we may assume $w_0=\bar w_0=1$, i.e.\ $w_N=\bar w_N=(-1)^N$ throughout their ranges.
Then in the one-bead sector, using $P_e+S_e=N-m_e$ and $P_f+S_f=N-m_f$ with
$N=\#(M\cap(b,b+g))$, Lemma~\ref{lem:onebead} gives
\[
  (1-2m_f)(-1)^{N-m_f}-(1-2m_e)(-1)^{N-m_e}=(-1)^N-(-1)^N=0,
\]
because $(1-2m)(-1)^{-m}=1$ for $m\in\{0,1\}$.  In the two-bead sector
(Lemma~\ref{lem:twobead} below) the element is
$\bar w_Q w_{P-k}-w_P\bar w_{Q+k}=(-1)^{P+Q-k}-(-1)^{P+Q+k}=0$.  By
Lemma~\ref{lem:sectors} these are all the matrix elements, so the commutator vanishes.

\emph{Necessity.}  Assume $[R_e^{w},R_f^{\bar w}]=0$, so \eqref{eq:B} holds.

\emph{Case 1: $w_0=0$ and $w\not\equiv0$.}  Let $i$ be least with $w_i\ne0$; then
$1\le i\le e-1$.  Taking $z=0$ in \eqref{eq:B} gives $w_u\bar w_v=0$ for all
$u\le e-1$, $v\le f-e-1$; with $u=i$ and $K$ a domain this forces
\begin{equation}\label{eq:c1a}
  \bar w_v=0\qquad (0\le v\le f-e-1).
\end{equation}
Now take $z=i$ in \eqref{eq:B}:  $w_i\bar w_{u+v+1}=-w_u\bar w_{v+i}$.
For $u<i$ the right side is $0$, so $\bar w_n=0$ for all $n$ of the form $u+v+1$ with
$u\le i-1$, $v\le f-e-1$, i.e.\ for $1\le n\le i+f-e-1$; together with \eqref{eq:c1a},
\begin{equation}\label{eq:c1b}
  \bar w_n=0\qquad(0\le n\le i+f-e-1).
\end{equation}
For $i\le u\le e-1$ the right side is $-w_u\bar w_{v+i}$ and $v+i\le i+f-e-1$, so
$\bar w_{v+i}=0$ by \eqref{eq:c1b}; hence $\bar w_{u+v+1}=0$ for
$i\le u\le e-1$, $0\le v\le f-e-1$, i.e.\ for $i+1\le n\le f-1$.  With \eqref{eq:c1b}
this gives $\bar w\equiv0$: case (b).

\emph{Case 2: $\bar w_0=0$ and $w_0\ne0$.}  Taking $z=0$, $v=0$ in \eqref{eq:B} gives
$w_0\bar w_{u+1}=-w_u\bar w_0=0$, so $\bar w_n=0$ for $1\le n\le e$; with $\bar w_0=0$,
$\bar w_n=0$ for $n\le e$.  Inductively, suppose $\bar w_n=0$ for all $n\le m$ with
$m\ge e$ and $m+1\le f-1$.  Write $m=u+v$ with $0\le u\le e-1$, $0\le v\le f-e-1$
(possible since $m\le f-2=(e-1)+(f-e-1)$).  Then $v\le m$, so \eqref{eq:B} with $z=0$
gives $w_0\bar w_{m+1}=-w_u\bar w_v=0$, whence $\bar w_{m+1}=0$.  So $\bar w\equiv0$:
case (b).

\emph{Case 3: $w_0\ne0$ and $\bar w_0\ne0$.}  Rescaling, assume $w_0=\bar w_0=1$.  Two
specialisations of \eqref{eq:B}:
\[
  \text{(B0)}\quad \bar w_{u+v+1}=-w_u\bar w_v\ \ (z=0),
  \qquad
  \text{(B1)}\quad w_z\bar w_{v+1}=-\bar w_{v+z}\ \ (u=0).
\]
From (B0) at $u=v=0$: $\bar w_1=-1$.  From (B1) at $v=0$: $w_z\bar w_1=-\bar w_z$, i.e.
\begin{equation}\label{eq:same}
  w_z=\bar w_z\qquad(0\le z\le e-1).
\end{equation}
From (B0) at $v=0$ and \eqref{eq:same}: $\bar w_{u+1}=-w_u=-\bar w_u$ for
$0\le u\le e-1$, so $\bar w_n=(-1)^n$ for $n\le e$.  From (B0) at $u=0$:
$\bar w_{v+1}=-\bar w_v$ for $0\le v\le f-e-1$, so $\bar w_n=(-1)^n$ for $n\le f-e$.
Let $n$ satisfy $\max(e,f-e)<n\le f-1$ and assume $\bar w_m=(-1)^m$ for all $m<n$.
Write $n-1=u+v$ with $0\le u\le e-1$, $0\le v\le f-e-1$ (possible since
$n-1\le f-2$).  Both $u\le e-1<n$ and $v\le f-e-1<n$, so by (B0) and \eqref{eq:same}
\[
  \bar w_n=-w_u\bar w_v=-\bar w_u\bar w_v=-(-1)^{u+v}=(-1)^{u+v+1}=(-1)^n .
\]
Hence $\bar w_n=(-1)^n$ for all $0\le n\le f-1$, and $w_n=\bar w_n=(-1)^n$ for
$0\le n\le e-1$ by \eqref{eq:same}: case (c) with $w_0=\bar w_0=1$.

Finally, in case (c) with $w_0=1$, $R_e^{w}s_\lambda=\sum_\mu(-1)^{\hgt(\mu/\lambda)}
s_\mu$, which is the Murnaghan--Nakayama rule for $p_e\,s_\lambda$
(Remark~\ref{rem:mn}).
\end{proof}

\begin{corollary}[answer to the primary question, in $\tau$ coordinates]\label{cor:tau}
Let $\tau_N=(-1)^Nw_N$ be the fermionic weights of the brief, normalised by $\tau_0=1$,
shared across $e$.  For any $e\ne f$, the locus
$\{\tau:[R_e^{\tau},R_f^{\tau}]=0\}$ is the single point $\tau_N=1$
$(0\le N\le\max(e,f)-1)$.  In particular for $(e,f)=(1,3)$ --- the first pair that can
distinguish multi-$t$ from single-$t$, since $R_3^{\tau}$ is the first operator carrying
two free parameters --- the locus in the $(\tau_1,\tau_2)$-plane is the point $(1,1)$.
It meets the single-$t$ parabola $\tau_2=\tau_1^2$ only there, and it contains nothing
else: there is no new commuting family off the single-$t$ curve.
\end{corollary}

\begin{proof}
Theorem~\ref{thm:main} case (c) with $w_0=\bar w_0=1$ reads $\tau_N=1$; cases (a), (b)
are excluded by $\tau_0=1$.  For $(e,f)=(1,3)$ relation \eqref{eq:B} has $u=z=0$ and
$v\in\{0,1\}$, giving $\bar w_1=-1$ and $\bar w_2=-\bar w_1=1$, i.e.\
$\tau_1=\tau_2=1$.
\end{proof}

\begin{corollary}[the family]\label{cor:family}
Let $E\subseteq\Z_{\ge1}$ with $|E|\ge2$ and let $R_e^{w^{(e)}}$, $e\in E$, be nonzero
deformed ribbon operators that pairwise commute.  Then for every $e\in E$,
$w^{(e)}_N=(-1)^Nw^{(e)}_0$ and $R_e^{w^{(e)}}=w^{(e)}_0\,\Mult{p_e}$.  In particular the
only commutative algebra generated by deformed ribbon operators of at least two distinct
sizes is the algebra of power-sum multiplications.
\end{corollary}

\section{Why: the two sectors cut out different loci}

Theorem~\ref{thm:main} used only the one-bead sector.  The two-bead sector gives an
independent second route to (most of) the conclusion, and comparing the two says exactly
where the extra freedom of a free weight dies.

\begin{lemma}[sectors]\label{lem:sectors}
$\langle s_\mu,[R_e^{w},R_f^{\bar w}]s_\lambda\rangle=0$ unless
$|M(\lambda)\,\triangle\,M(\mu)|\in\{2,4\}$.
\end{lemma}

\begin{proof}
Each $R^{\bullet}_g$ moves one bead, so a product moves at most two; and $M(\mu)=M(\lambda)$
is impossible since $|\mu|=|\lambda|+e+f>|\lambda|$.
\end{proof}

\begin{lemma}[two-bead element]\label{lem:twobead}
Let $e\ne f$ and $M(\mu)=(M\setminus\{b,c\})\cup\{b+e,c+f\}$ with $b,c\in M$,
$b+e,c+f\notin M$ and $b,c,b+e,c+f$ pairwise distinct.  Put $P=\#(M\cap(b,b+e))$,
$Q=\#(M\cap(c,c+f))$ and $k=[\,b+e\in(c,c+f)\,]-[\,b\in(c,c+f)\,]\in\{-1,0,1\}$.  Then
\[
  \bigl\langle s_\mu,[R_e^{w},R_f^{\bar w}]s_\lambda\bigr\rangle
  \;=\;\bar w_{Q}\,w_{P-k}\;-\;w_{P}\,\bar w_{Q+k}.
\]
Every other $|M\triangle M(\mu)|=4$ element vanishes.
\end{lemma}

\begin{proof}
This is \cite[Thm.~5(2)]{Clio0907} with the weights left free; the route enumeration is
unchanged.  No cancellation occurs, so the removed sites are $M\setminus M(\mu)$ and the
added sites are $M(\mu)\setminus M$; for $e\ne f$ the pair $(b,c)$ is determined ($b$ is
the unique removed site with $b+e$ added).  Performing the $e$-move first gives heights
$P$ and $Q+k$ (the $e$-move changes the count in $(c,c+f)$ by $k$); performing the
$f$-move first gives $Q$ and $P-k$ (by \cite[Lem.~4]{Clio0907}, the two interference
indices sum to zero off the diagonals $b=c$, $b+e=c+f$).
\end{proof}

\begin{proposition}[the two-bead sector restores a single parameter, and only that]
\label{prop:twobead}
Let $2\le e<f$ and $w_0=\bar w_0=1$.  If every two-bead matrix element of
$[R_e^{w},R_f^{\bar w}]$ vanishes, then there is $s\in K^{\times}$ with
\[
   w_N=s^{N}\ (N\le e-1),\qquad \bar w_N=s^{-N}\ (N\le f-1).
\]
If in addition the weights are shared ($w=\bar w$ on the common range and $e\ge2$), then
$s^2=1$.  Conversely every such pair kills the two-bead sector.
\end{proposition}

\begin{proof}
Take $b=0$ and $c=d$ with $1\le d\le e-1$; then $b<c<b+e<c+f$ (the last because
$d\ge1>e-f$), so the four sites are distinct, $b+e\in(c,c+f)$, $b\notin(c,c+f)$ and
$k=1$.  Decompose $(0,e)=(0,d)\sqcup\{d\}\sqcup(d,e)$ and
$(d,d+f)=(d,e)\sqcup\{e\}\sqcup(e,d+f)$.  Prescribing (Lemma~\ref{lem:realise}) $a$
beads in $(0,d)$, $m$ beads in $(d,e)$, $q$ beads in $(e,d+f)$ and none at $e$ or $d+f$,
we get $P=a+1+m$ and $Q=m+q$, with $a\le d-1$, $m\le e-d-1$, $q\le d+f-e-1$.  Choosing
$d=e-1$ forces $m=0$ and leaves $a\le e-2$, $q\le f-2$ free and independent.  The
vanishing condition $\bar w_Qw_{P-1}=w_P\bar w_{Q+1}$ becomes
\[
   \bar w_{q}\,w_{a}\;=\;w_{a+1}\,\bar w_{q+1}\qquad(0\le a\le e-2,\ 0\le q\le f-2).
\]
At $a=q=0$: $1=w_1\bar w_1$, so $s:=w_1$ is a unit with $\bar w_1=s^{-1}$.  At $a=0$:
$\bar w_q=s\,\bar w_{q+1}$, so $\bar w_n=s^{-n}$ for $n\le f-1$.  At $q=0$:
$w_a=w_{a+1}s^{-1}$, so $w_n=s^{n}$ for $n\le e-1$.  If the weights are shared then
$s=w_1=\bar w_1=s^{-1}$.  The converse is the computation
$\bar w_Qw_{P-k}-w_P\bar w_{Q+k}=s^{-Q+P-k}-s^{P-Q-k}=0$.
\end{proof}

\begin{remark}[the shape of the answer]\label{rem:shape}
Proposition~\ref{prop:twobead} is the precise sense in which the multi-$t$ freedom is an
illusion, and it is worth separating from Theorem~\ref{thm:main}:
\begin{itemize}
\item The \textbf{two-bead sector} forces both weights to be \emph{multiplicative}:
      $w_N=s^N$, $\bar w_N=s^{-N}$.  That is precisely the statement
      $\tau_2=\tau_1^{2}$ and its higher analogues --- a single parameter, restored.  It
      does \emph{not} determine $s$.
\item The \textbf{one-bead sector} forces the anchor outright, with no multiplicativity
      input at all (Theorem~\ref{thm:main}).
\end{itemize}
So the two mechanisms are genuinely different, and the one-bead mechanism is the stronger
of the two.  Note also that the two-bead sector is empty when $\min(e,f)=1$
(\cite[Cor.~9]{Clio0907}: $k\ne0$ forces one of the four sites to collide), which is why
the theorem must rest on the one-bead sector: the pair $(1,3)$ --- the cheapest genuine
test --- has no two-bead sector at all.
\end{remark}

\begin{remark}[the mechanism, stated once more without indices]\label{rem:why}
Read \eqref{eq:B} at $u=z=0$: $\bar w_{v+1}=-\bar w_v$.  On the left, the site $b+e$ is a
bead \emph{jumped over} by the long $f$-move $b\to b+f$, and so is counted by its height.
On the right, the same site is the \emph{hinge} at which the $(e+f)$-ribbon factors, and
a hinge is an endpoint of both open intervals, hence counted by neither height.  One
occupied site, two readings, differing by exactly $1$.  Closure of the bracket demands
that the weight convert that unit shift into a sign, and the only weights that can do so
are the alternating ones.  The scalar $1+t$ of the single-parameter theorem is the
instance $v=0$ of this relation --- it was never a fact about $t$, and giving $t$ more
room to move was never going to help.
\end{remark}

\section{Computational verification}\label{sec:ver}

All code is in \texttt{probes/2026-09-11-c2-Q147/}.  The engine
(\texttt{engine.py}) was written fresh for this note; in particular it never forms
$(-t)^N$ or $t^{\hgt}$, and the weight enters only as an opaque symbol.

\begin{enumerate}
\item \textbf{Two independent ribbon enumerations agree.}  \texttt{ribbons\_young}
  enumerates all $\mu\vdash|\lambda|+e$ containing $\lambda$ and tests the skew shape for
  the $2\times2$ and connectivity conditions, computing the height as (rows)$-1$;
  \texttt{ribbons\_maya} performs bead moves $b\to b+e$ and computes the height as beads
  jumped.  They agree as multisets of (shape, height) for every $\lambda$ with
  $|\lambda|\le8$ and every $e\le5$.

\item \textbf{Lemma~\ref{lem:onebead} verified.}  For
  $(e,f)\in\{(1,2),(1,3),(2,3),(2,4),(3,4),(2,5),(3,5)\}$, all four occupancy patterns
  $(m_e,m_f)$ and every bead pattern in the three open sub-intervals --- $316$
  configurations, built by the construction of Lemma~\ref{lem:realise} --- the formula
  agrees with the engine's commutator on the nose, with \emph{independent} symbolic
  weights.  $0$ mismatches.

\item \textbf{Lemma~\ref{lem:twobead} verified} likewise on $42$ configurations
  ($k=+1$, $(e,f)\in\{(2,3),(2,4),(3,4),(3,5)\}$).  $0$ mismatches.

\item \textbf{The locus, computed without the theorem.}  Collecting every nonzero entry
  of $[R_e^{w},R_f^{\bar w}]s_\lambda$ over all $\lambda$ up to size $13$ and solving the
  resulting polynomial system:
  \begin{center}
  \begin{tabular}{lll}
  \hline
  $(e,f)$ & shared weights & independent weights\\
  \hline
  $(1,2)$ & $w_1=-1$ & --- \\
  $(1,3)$ & $w_1=-1,w_2=1$ & $\bar w_1=-1,\bar w_2=1$\\
  $(1,4)$, $(1,5)$ & $w_N=(-1)^N$, $N\le f-1$ & --- \\
  $(2,3)$, $(2,4)$, $(2,5)$, $(2,6)$ & $w_N=(-1)^N$, $N\le f-1$ &
      $w_N=\bar w_N=(-1)^N$\\
  $(3,4)$, $(3,5)$, $(3,6)$, $(4,5)$ & $w_N=(-1)^N$, $N\le f-1$ &
      $w_N=\bar w_N=(-1)^N$\\
  \hline
  \end{tabular}
  \end{center}
  In every case the solution set is the single point of Theorem~\ref{thm:main}(c), and
  every symbol visible to the pair appears in the equations.

\item \textbf{The sector split verified.}  Separating the equations by
  $|M\triangle M'|$: the one-bead equations alone have the unique solution
  $w_N=\bar w_N=(-1)^N$; the two-bead equations alone have solution set
  $\{w_N=s^N,\ \bar w_N=s^{-N}\}$ for independent weights and
  $\{w_N=1\}\cup\{w_N=(-1)^N\}$ for shared weights.  Checked for
  $(e,f)\in\{(2,3),(2,4),(3,4),(2,5),(3,5)\}$.  This is
  Proposition~\ref{prop:twobead} and Remark~\ref{rem:shape}.

\item \textbf{The anchor identification, independently.}  That $w_N=(-1)^N$ gives
  multiplication by $p_e$ was checked without any ribbon, abacus or
  Murnaghan--Nakayama input: $s_\lambda$ was expanded in the monomial basis by counting
  semistandard tableaux (Kostka numbers), multiplied by $p_e=\sum_ix_i^e$ in
  $n=|\lambda|+e$ variables, and re-expanded in the Schur basis by inverting the
  (unitriangular) Kostka matrix.  Agreement on $76/76$ pairs with $e\le4$,
  $|\lambda|\le5$.

\item \textbf{The degenerate stratum.}  Dropping the normalisation $w_0=1$ entirely and
  solving, the locus decomposes exactly as Theorem~\ref{thm:main}(a)--(c), with no
  further component: verified for $(e,f)\in\{(2,3),(2,4),(3,4)\}$.
\end{enumerate}

\subsection*{Negative controls}

\begin{enumerate}
\item[(N1)] \textbf{The solver manufactures nothing.}  Run on $e=f=3$, where the
  commutator is identically zero, the pipeline returns $0$ equations and hence no
  constraint.  (A pipeline that produced constraints here would be producing them from
  its own bookkeeping.)
\item[(N2)] \textbf{The test is not degenerate in the new direction.}  This is the
  control that matters, because the whole question is whether $\tau_2$ is genuinely
  probed.  For $(e,f)=(1,3)$ the equation set is $\{\pm(w_1+1),\pm(w_1+w_2)\}$;
  substituting the forced value $w_1=-1$ leaves the residual equation $w_2-1$, which
  \emph{does} contain $w_2$.  So $w_2$ is constrained by a relation that survives the
  determination of $w_1$: the vanishing of the new freedom is measured, not assumed.
\item[(N3)] \textbf{The deformation is real.}  $R_3^{w}s_{(1)}=s_{(4)}+w_1s_{(2,2)}
  +w_2s_{(1,1,1,1)}$: the three weights sit on three distinct targets, so $w_2$ and
  $w_1^2$ are separated by the operator itself.  A deformation invisible to the operator
  would make items 4 and 5 vacuous.
\item[(N4)] \textbf{$[R_1,R_2]$ is excluded by construction.}  It is nowhere used;
  it sees only $w_1$ and so cannot distinguish free weights from a single parameter.
\end{enumerate}

\section{What is and is not proved}\label{sec:gaps}

\textbf{Proved.}  Theorem~\ref{thm:main} (all cases, all $1\le e<f$, independent weights,
any integral domain), Corollaries~\ref{cor:tau} and~\ref{cor:family},
Proposition~\ref{prop:twobead}.  Lemmas~\ref{lem:realise}, \ref{lem:onebead},
\ref{lem:sectors} are proved here; Lemma~\ref{lem:twobead} is the free-weight form of
\cite[Thm.~5(2)]{Clio0907} and is used only for Proposition~\ref{prop:twobead} and for
the sufficiency half of Theorem~\ref{thm:main}, where it is applied at the explicit
weights $w_N=(-1)^N$.

\textbf{Not proved / out of scope.}
\begin{enumerate}
\item \emph{Weights depending on more than the height.}  A weight $w_{N}$ could be
  replaced by a weight depending on the whole occupancy word of the interval
  $(b,b+e)$, not only on the number of beads in it.  That is a strictly larger
  deformation ($2^{e-1}$ parameters rather than $e$) and this note says nothing about it.
  The mechanism of Remark~\ref{rem:why} suggests it will not help --- the hinge/jumped
  dichotomy is a statement about one site --- but the bookkeeping is different and I have
  not done it.
\item \emph{$q$-deformed Clifford relations.}  The deformation here is diagonal: the
  fermions stay undeformed and only the dressing moves.  Deforming
  $\{\psi_a,\psi_b^{*}\}$ is a different question.
\item \emph{Level $\ell>1$.}  Everything above is level $1$.
\item \emph{The adjoint relation.}  I did not redo $[R_e^{*},R_e]=[e]_{t^2}$ in the
  $\tau$ variables, as the brief permitted but did not require; it is not used anywhere
  here.
\item \emph{$e=f$.}  Trivially $[R_e^w,R_e^w]=0$ for every $w$; the theorem is about
  distinct sizes only, and this is not a gap but the hypothesis.
\end{enumerate}

\textbf{Effect on the earlier theorem.}  Nothing is retracted.  The single-parameter
forcing theorem --- $R_e(t)$ and $R_f(t)$ do not commute for $e\ne f$, and $(1+t)$
divides every matrix element \cite[Cor.~11]{Clio0906}, sharply \cite[Cor.~7]{Clio0907}
--- is the restriction of Theorem~\ref{thm:main} to the curve $w_N=t^N$, and remains
correct: there case~(c) with $w_0=1$ reads $t^N=(-1)^N$, i.e.\ $t=-1$.  What changes is its standing: the hypothesis
``one parameter $t$'' was not load-bearing, and the conclusion ``$t=-1$'' was a shadow of
the sharper statement ``the height weight must alternate''.

\begin{remark}\label{rem:mn}
$R_e^{((-1)^\bullet)}=\Mult{p_e}$ is the Murnaghan--Nakayama rule
$p_e\,s_\lambda=\sum_\mu(-1)^{\hgt(\mu/\lambda)}s_\mu$, verified independently in
\S\ref{sec:ver}(6).  The proof of Theorem~\ref{thm:main} does not use it: sufficiency is
checked directly on matrix elements, so the theorem is self-contained and the
identification is a reading of the answer, not a step in it.
\end{remark}

\begin{thebibliography}{9}
\bibitem{Clio0906} Clio, \emph{A fermionic normal form for the ribbon operator $R_e(t)$,
  and its identification with the Kashiwara--Miwa--Stern boson}, 6 September 2026,
  \texttt{proofs/2026-09-06-c2-Q91-fermionic-normal-form.tex}.
\bibitem{Clio0907} Clio, \emph{The cross-rank ribbon commutator}, 7 September 2026,
  \texttt{proofs/2026-09-07-Q92-cross-rank-commutator.tex}.
\bibitem{LLT} A.~Lascoux, B.~Leclerc, J.-Y.~Thibon, \emph{Green polynomials and
  Hall--Littlewood functions at roots of unity}, Lett.\ Math.\ Phys.\ \textbf{35} (1995)
  359--374.  (The move of replacing a parameter raised to a statistic by a parameter
  indexed by it.)
\end{thebibliography}

\end{document}
